% page 167 du fac-simile (image p-166 du scan)
% bloc 170.2 x 259.3 mm ; marge gauche 31.8 mm
\pgc{10.160mm}{0.000mm}{
\l{}
\l{}
\l{\cel{54}157.}
\l{}
\l{}
\l{falco. Entaliuro en la kadro di pordo, di fenestro, por ke la}
\l{\cel{3}bordo di la klapo qua klozas lu esez kontre-aplikebla sen}
\l{\cel{3}lasar lum-lakuno. - DR.}
\l{}
\l{falchar. (\sou{trans.})\cel{2}(\sou{agrokultivo}). Tranchar (la feno, la cereali)}
\l{\cel{3}per instrumento ek longa mancho ligna e granda lamo arkatra.-FI.}
\l{}
\l{faldar. (\sou{trans.}) Igar duopla unfoye o plurafoye (materio flexebla,}
\l{\cel{3}stofo, papero, e c.) per abatar ica parto sur ita. - DE.}
\l{}
\l{faleno. (\sou{zool.}) Nokto-papiliono. - L.\cel{1}\sou{phanaena}. - DEFIS.}
\l{}
\l{faliar.\cel{2}I. (\sou{trans.}) Ne atingar to quon onu vizis. - II. (\sou{netrans.})}
\l{\cel{3}Ne agar to quon onu devas etikale. - III. (\sou{netrans.}) (\sou{komerco)}}
\l{\cel{3}Deklarar su nesolventa, cesar la pagado di sua debaji. - DEFIS.}
\l{}
\l{falko. (\sou{navo}). Plankaro extera per qua onu plualtigas la bordo di}
\l{\cel{3}navo kande la ondi esas tro granda. - FIS.}
\l{}
\l{falkono. (zool.)\cel{2}Raptucelo - genero di la klaso \textquotedbl{}rapteri\textquotedbl{} - quan}
\l{\cel{3}onu dresis olim por chasado. - L.\cel{1}\sou{falco}. - DEFIS.}
\l{}
\l{falsa. I. Qua havas, di kozo, nur la semblo. - II. Qua prenas, por}
\l{\cel{3}trompar, la semblo di kozo. - DEFIRS.}
\l{}
\l{falseo. Voco akuta, produktita da la vibro di la kordi supra di}
\l{\cel{3}la laringo. - DEFIRS.}
\l{}
\l{\sou{falva}.\cel{2}Flava rufe. - DF.}
\l{}
\l{famo. Favoreso, segun-modeso, quan la nomo di ulu (o di ulo)}
\l{\cel{3}obtenis. - DEFI.}
\l{}
\l{familio. I. Asemblitaro de la personi quin unionas la sango o}
\l{\cel{3}la alianco, qui vivas en un hemo, partikulare la patrulo, la}
\l{\cel{3}patrino, e lia filii. -\cel{2}II. Asemblitaro de la personi unionita}
\l{\cel{3}per la sango o la alianco. - III. La serio de la personi qui}
\l{\cel{3}devenis de la sama familio, generacionope. IV. Asemblitaro de}
\l{\cel{3}enti di qui la origino esas komuna, di qui la interesti esas}
\l{\cel{3}komuna. - V.(\sou{zool. e bot}.) Kolekturo de generi intervicina pro}
\l{\cel{3}ula karakterizivi generala quin li posedas komune. - DEFIRS.}
\l{}
\l{\sou{familiara}. I. Qua relatas ad ulu libere, sen jeneso, sen ceremo-}
\l{\cel{3}nialo. - II. Qua kustumas la uzo di ulo ; di qua la praktiko}
\l{\cel{3}la uzo, esas ordinara ye ulu.- DEFIS.}
\l{}
\l{\sou{famino.}\cel{2}Sufrado generala, produktita da la indijo di nutrivi.EFI.}
\l{}
\l{\sou{fanar}. (trans.) Netigar, purigar (frumento) per sukusar lu sur}
\l{\cel{3}sorto di korbo plata, di qua la bordo esas basa, por forjetar}
\l{\cel{3}la polvo, la brano, la palio-peceti. - DEF.}
\l{}
\l{\sou{fanatika}. I.Qua kredas inspiresar deale. - II. Ta quan zelo reli-}
\l{\cel{3}giala blindigas, incitas ad agi neyusta, violentoza, qui trans-}
\l{\cel{3}\sou{iras}\cel{1}la limiti di la raciono komuna. - DEFIRS.}
}
